\section{十三、Wi-Fi 与蓝牙：从内存帧到无线电波}

有线以太网把 MAC 帧交给 PCS、PMA 和 PMD；无线接口还要在 MAC 与天线之间完成信道接入、基带调制、数模转换、射频变频和功率放大。CPU 看到的仍是描述符、DMA 和完成项，但“介质正在被谁使用”不再由独占线对决定，而由共享频谱上的侦听、退避、确认和重传共同决定。

\begin{pdfdiagram}[x=1cm,y=1cm]
  \node[hw memory, minimum width=2.3cm] (mem) at (1.3,3.7) {主存\\帧/描述符};
  \node[hw control, minimum width=2.3cm] (mac) at (4.1,3.7) {MAC\\排队/退避/重传};
  \node[hw compute, minimum width=2.3cm] (bb) at (6.9,3.7) {基带\\编码/调制/OFDM};
  \node[hw compute, minimum width=2.3cm] (rf) at (9.7,3.7) {RF 收发机\\DAC/ADC/变频};
  \node[hw block, minimum width=2.3cm] (ant) at (12.5,3.7) {前端与天线\\PA/LNA/滤波};
  \draw[hw bus] (mem) -- (mac);
  \draw[hw bus] (mac) -- (bb);
  \draw[hw bus] (bb) -- (rf);
  \draw[hw bus] (rf) -- (ant);

  \node[hw state, minimum width=2.5cm] (air) at (12.5,1.0) {共享无线信道\\噪声/干扰/多径};
  \draw[hw bus] (ant) -- (air);
  \node[hw control, minimum width=2.8cm] (ctrl) at (8.3,1.0) {信道估计与速率控制\\RSSI/SNR/EVM/重传率};
  \draw[hw return] (air) -- (ctrl);
  \draw[hw return] (ctrl) -- (mac);
\end{pdfdiagram}
\diagramnote{无线发送与反馈是一个闭环。上排把内存帧逐层转换为射频能量，下排把信道测量、确认与重传结果送回速率控制和 MAC。吞吐下降可能来自队列、争用、基带误码、射频功率或天线环境，必须按层读取证据。}

\topic{发送路径把 bit 映射到频率、相位和幅度}

驱动写好发送描述符后，设备 DMA 读取帧和控制信息。MAC 添加序号、地址、校验和无线控制字段，并等待信道接入。基带先执行扰码、前向纠错编码和交织，再把 bit 组映射成调制符号；OFDM 把符号分配到许多正交子载波，经 IFFT 形成时域采样并加入循环前缀。DAC 把数字采样变成模拟基带或中频信号，混频器上变频到射频，功率放大器经滤波与天线把能量送入空间。

高阶调制让每个符号携带更多 bit，却要求更高信噪比和更低误差向量幅度。速率控制根据接收信号强度、信噪比、重传率和历史成功概率选择调制编码方案、信道宽度与空间流数量。最大标称速率只在足够信道质量、足够空间流和较少竞争时成立。

\topic{接收路径先恢复模拟波形，再判断 MAC 帧}

天线接收的信号先经带通滤波和低噪声放大器，再由混频器下变频；ADC 产生数字采样，自动增益控制避免过小或饱和。基带利用前导序列完成包检测、频偏估计、定时同步和信道估计，再执行 FFT、均衡、解调、解交织和纠错译码。只有通过物理层与 MAC 校验的帧，才被写入接收缓冲并产生完成项。

接收描述符耗尽时，即使射频和基带已正确解出帧，设备也没有合法 DMA 目标；相反，FCS、PHY 或重试计数上升时，增加软件接收环深度并不能修复信道问题。定位必须同时读取 RX ring 水位、MAC 丢弃、PHY 错误、RSSI/SNR、调制方案和重传计数。

\topic{关联与密钥建立发生在普通数据传输之前}

Wi-Fi 站点先扫描信道并发现网络，选择接入点后执行认证、关联和密钥协商，再安装用于数据帧保护的密钥。关联成功只证明控制状态已经建立；DHCP、路由和应用连接仍属于后续软件路径。漫游时，新接入点、密钥和队列代际必须在旧链路失效前后明确切换，避免旧重传帧被新连接错误接受。

蓝牙使用相同的“基带 + RF + 天线”物理骨架，但面向低功耗短包和跳频信道。控制器按连接事件唤醒，在指定频点交换少量包，然后重新休眠。低功耗来自缩短射频开启时间，不是让无线状态完全消失；时钟偏差、连接间隔、监督超时和重传仍决定连接能否持续。

\begin{longtable}{L{0.18\linewidth}L{0.31\linewidth}L{0.39\linewidth}}
\toprule
证据 & 更可能对应的边界 & 不应直接采取的动作 \\
\midrule
TX ring 长期积压 & MAC 争用、固件停顿或链路未关联 & 直接提高发射功率 \\\tablerowrule
重传率与 FCS 错误同时升高 & 干扰、多径、弱信号或频偏 & 只增加驱动队列深度 \\\tablerowrule
RSSI 高但吞吐低 & 同信道竞争、低 SNR、接收饱和或队列限制 & 仅依据信号格数判定链路良好 \\\tablerowrule
PHY 成功但主机无包 & RX 描述符、DMA、IOMMU 或中断路径 & 重新调天线匹配 \\\tablerowrule
关联反复重置 & 密钥、漫游、固件状态或时钟/电源问题 & 把每次断开都解释为射频弱 \\
\bottomrule
\end{longtable}

\section{十四、RDMA、RoCE 与 SmartNIC：网络设备直接访问应用缓冲区}

普通网络接收通常经过内核协议栈和 socket 缓冲，再复制或映射给应用。RDMA 把应用注册的内存区域、队列对和完成队列直接落实为 NIC 可访问的硬件对象，使远端请求可以把数据放入目标缓冲区，CPU 只处理建立、权限和完成。减少 CPU 参与并没有减少状态；它把更多所有权、顺序和错误恢复状态移动到了 RNIC。

\begin{pdfdiagram}[x=1cm,y=1cm]
  \node[hw state, minimum width=2.4cm] (app) at (1.3,3.4) {应用\\工作请求 WR};
  \node[hw memory, minimum width=2.4cm] (qp) at (4.2,3.4) {发送/接收队列\\QP 与 WQE};
  \node[hw control, minimum width=2.7cm] (rnic) at (7.3,3.4) {RNIC/SmartNIC\\传输、重传、权限};
  \node[hw block, minimum width=2.5cm] (wire) at (10.5,3.4) {Ethernet/RoCE\\或 InfiniBand};
  \draw[hw bus] (app) -- (qp);
  \draw[hw bus] (qp) -- (rnic);
  \draw[hw bus] (rnic) -- (wire);

  \node[hw memory, minimum width=2.7cm] (remote) at (10.5,1.0) {远端已注册内存\\rkey + 地址范围};
  \node[hw control, minimum width=2.7cm] (rrnic) at (7.3,1.0) {远端 RNIC\\IOMMU/页固定};
  \node[hw state, minimum width=2.4cm] (cq) at (4.2,1.0) {完成队列 CQ\\状态与代际};
  \node[hw state, minimum width=2.4cm] (done) at (1.3,1.0) {应用观察完成};
  \draw[hw bus] (wire) -- (remote);
  \draw[hw return] (remote) -- (rrnic);
  \draw[hw return] (rrnic) -- (cq);
  \draw[hw return] (cq) -- (done);
\end{pdfdiagram}
\diagramnote{一次 RDMA 操作的请求与完成环。上排由应用发布工作请求并经 RNIC 发送，下排由远端 RNIC 验证 rkey 和地址范围、访问已注册内存，再把完成返回本地 CQ。所谓“零拷贝”只减少 CPU 数据搬运，页固定、权限、队列槽、重传和完成代际仍必须闭合。}

\topic{内存注册把虚拟缓冲区变成设备可验证的能力}

应用注册内存时，内核固定或按需管理相关页，建立 IOMMU 映射，并向 RNIC 写入地址范围、权限和密钥。设备访问请求同时携带远端虚拟地址与 rkey；RNIC 查内存保护表，只有范围与权限都匹配才发起 DMA。rkey 不是加密密钥，而是不可猜测的能力标识；解除注册时必须阻止新请求、耗尽在途访问、撤销设备缓存的翻译和密钥，最后才能复用页框。

\topic{队列对与完成队列保存端到端顺序}

应用把 SEND、RECV、RDMA READ、RDMA WRITE 或原子操作写成 WQE，发布屏障后更新生产者指针。RNIC 读取 WQE、分段成网络包并维护包序、确认和重传。完成队列项表示该工作请求达到协议规定的本地完成边界，但不同操作的远端可见语义不同：本地 RDMA WRITE 完成通常说明数据已由传输层接受并抵达目标 RNIC/内存路径，不自动表示远端应用已经处理；需要应用级确认时仍要发送独立消息或更新双方约定的同步位置。

RoCE 把 RDMA 传输映射到以太网/IP 网络。无损网络可使用优先级流控与拥塞通知，但暂停帧会把一个拥塞点扩散到其他流量，拥塞控制仍要限制注入速率。丢包、ECN、重传超时、QP 状态和交换机队列应进入同一时间线，否则容易把网络拥塞误判为远端 CPU 延迟。

\topic{SmartNIC 与 DPU 把更多控制放到网卡侧}

SmartNIC 可以在网卡上执行虚拟交换、隧道封装、校验、加密和存储协议处理；DPU 进一步集成通用或专用处理核、内存与管理面。卸载后的数据路径可以绕过主机内核，但控制面仍要分配队列、IOMMU 域、密钥、带宽和复位代际。网卡固件异常时，主机必须能够隔离 DMA、收集设备日志并执行 FLR 或管理控制器复位。

\topic{PTP 硬件时间戳把“事件发生”固定在 MAC/PHY 边界}

软件在发送系统调用前后读取时钟，测到的是调度、队列和协议栈混合延迟。PTP 硬件时间戳在帧真正进入或离开 MAC/PHY 的固定位置锁存设备时钟，把可变排队部分排除。驱动把时间戳与描述符代际匹配后交给时间同步程序，程序估计时钟偏移与链路时延，再通过 PHC 或系统时钟调节接口校正频率和相位。时间戳若与错误队列项配对，偏差会表现为偶发的大幅跳变。

\section{十五、CAN 与 LIN：共享线上的确定性控制与错误恢复}

CAN 和 LIN 面向车辆与工业控制。它们的数据率低于 PCIe 和 Ethernet，却把电气鲁棒性、优先级仲裁、错误隔离和最坏响应时间放在首位。理解这两类总线可以看到：协议带宽较低并不意味着控制状态较少，反而常要求更明确的故障升级规则。

\begin{pdfdiagram}[x=1cm,y=1cm]
  \node[hw state, minimum width=2.4cm] (n0) at (1.6,4.1) {节点 A\\ID 0x120};
  \node[hw state, minimum width=2.4cm] (n1) at (4.7,4.1) {节点 B\\ID 0x080};
  \node[hw state, minimum width=2.4cm] (n2) at (7.8,4.1) {节点 C\\ID 0x300};
  \draw[BookMuted, line width=1.1pt] (1.6,2.65) -- (9.3,2.65);
  \draw[hw bus] (n0.south) -- (1.6,2.65);
  \draw[hw bus] (n1.south) -- (4.7,2.65);
  \draw[hw bus] (n2.south) -- (7.8,2.65);
  \node[hw block, minimum width=2.8cm] (bus) at (11.2,2.65) {CAN 差分总线\\显性 0 / 隐性 1};
  \draw[hw bus] (9.3,2.65) -- (bus);

  \node[hw control, minimum width=3.0cm] (arb) at (2.7,0.3) {逐 bit 仲裁\\发送 1 却读到 0 则退出};
  \node[hw control, minimum width=3.0cm] (err) at (6.8,0.3) {错误计数\\active/passive/bus-off};
  \node[hw state, minimum width=2.8cm] (winner) at (11.2,0.3) {最低 ID 继续发送\\其余节点等待重试};
  \draw[hw return] (bus.south) -- (11.2,1.45) -- (2.7,1.45) -- (arb.north);
  \draw[hw bus] (arb) -- (err);
  \draw[hw bus] (err) -- (winner);
\end{pdfdiagram}
\diagramnote{CAN 仲裁与错误隔离共享同一物理反馈：节点边发送边读取总线。仲裁阶段，读回更强的显性 0 表示自己失去优先级但不算错误；数据阶段，发送值与读回值不符会进入错误计数。协议根据当前阶段解释同一电气事实。}

\topic{CAN 用线与关系完成无破坏仲裁}

CAN 收发器把逻辑状态转换为 CAN\_H/CAN\_L 差分电压。显性 bit 0 可以覆盖隐性 bit 1，多个节点同时发送时，总线呈现逐 bit 线与结果。标识符从高位开始发送；某节点发送隐性 1 却读到显性 0，就知道另一个更低标识符帧拥有更高优先级，于是停止发送而不破坏获胜帧。最坏等待时间由高优先级帧频率、帧长度和总线利用率决定，不能只用平均带宽判断实时性。

数据阶段仍执行逐 bit 回读、位填充、CRC、格式和确认检查。节点发现错误后发送错误标志，使所有参与者丢弃该帧并重新发送。发送与接收错误计数器决定节点处于 error-active、error-passive 或 bus-off；故障节点若持续干扰总线，最终会被协议隔离。恢复 bus-off 必须满足平台规定的静默时间与管理策略，不能在电气短路仍存在时立即循环重启。

\topic{LIN 用主节点调度表换取更低成本}

LIN 通常使用单线加地、较低数据率和主从调度。主节点按调度表发送 break、同步字节和标识符，被指定的发布者随后发送数据与校验。没有逐 bit 多主仲裁，时间确定性来自主节点预先安排的时隙。节点可使用低成本 RC 时钟，通过同步字节校准当前帧的 bit 时间。

LIN 的低成本也带来边界：主节点故障会影响整条调度，单线对短路和地偏移更敏感，错误检测能力弱于 CAN。车辆常用 CAN 承担制动、动力等高可靠网络，用 LIN 承担车窗、座椅和传感器等低速支线；网关在两者之间转换信号时，必须定义更新时间、超时默认值和失联后的安全状态。
